问题一的核心不仅在于验证"鱼类资源量增加"，更在于同时揭示"鱼类增长的驱动机制"与"基础资源的承载压力变化"。若将四大家鱼合并为单一消费者功能群，将水草、浮游植物、浮游动物和底栖饵料简化为单一资源总量，虽仍能匹配2025年资源量恢复至1.8倍的观测事实，但无法揭示哪一层级资源率先承压，也难以解释题目中"水下荒漠"等底层生态退化现象。因此，模型必须保留四类基础资源与四大家鱼的基本食性对应关系。
